%% kapa: exam=1 round=34 year=2023 period=1 form=A subject=civil_law subject_ko=민법 seq=12 number=12
%% source: Q-Net 감정평가사 기출문제 (cst00309) artlSeq=5211789 file=raw/1008/5211789/2023년도 제34회 감정평가사 1차시험 1교시.pdf
%% conversion: kordoc 4.24.51 -> tools/gen1.py (auto); review=proofread-sonnet; flags=none
\begin{question}{12}
착오로 인한 의사표시에 관한 설명으로 옳지 않은 것은? (다툼이 있으면 판례에 따름)
\begin{choices}
  \item 매도인의 하자담보책임이 성립하더라도 착오를 이유로 한 매수인의 취소권은 배제되지 않는다.
  \item 계약 당시를 기준으로 하여 장래의 미필적 사실의 발생에 대한 기대나 예상이 빗나간 경우, 착오취소는 인정되지 않는다.
  \item 동기의 착오는 동기가 표시되어 해석상 법률행위의 내용으로 된 경우에 한해서만 유일하게 고려된다.
  \item 매매계약에서 매수인이 목적물의 시가에 관해 착오를 하였더라도 이는 원칙적으로 중요부분의 착오에 해당하지 않는다.
  \item 상대방이 표의자의 착오를 알면서 이용하였다면 표의자의 착오에 중대한 과실이 있더라도 착오취소가 인정된다.
\end{choices}
\end{question}
\answer{3}
